\documentclass[conference]{IEEEtran}
\IEEEoverridecommandlockouts
\usepackage{cite}
\usepackage{amsmath,amssymb,amsfonts}
\usepackage{graphicx}
\usepackage{textcomp}
\usepackage{xcolor}
\usepackage{booktabs}
\usepackage{multirow}
\usepackage{array}
\usepackage{url}

\begin{document}

\title{Private Collaborative Computing: An Experimental Performance and Crossover Evaluation of Threshold Fully Homomorphic Encryption and Multi-Party Computation}

\author{\IEEEauthorblockN{Collaborative Cryptography and Distributed Systems Research Group}
\IEEEauthorblockA{\textit{Department of Computer Science and Cybersecurity} \\
Reproducible Dockerized Experimental Benchmark Report}}

\maketitle

\begin{abstract}
Organizations frequently possess complementary sensitive financial datasets that could yield high-value analytical insights if processed jointly, yet regulatory mandates and competitive risk prohibit centralizing raw records. Fully Homomorphic Encryption (FHE) and Secure Multi-Party Computation (SMPC) offer two fundamentally different cryptographic paradigms for privacy-preserving collaborative computation: FHE achieves non-interactive $O(1)$-round evaluation over lattice ciphertexts at the expense of substantial polynomial arithmetic and ciphertext expansion, whereas secret-sharing SMPC achieves lightweight field arithmetic at the cost of interactive communication rounds and Beaver multiplication triple preprocessing. In this paper, we present a rigorous empirical performance and crossover evaluation comparing a 3-party threshold Brakerski/Fan-Vercauteren (BFV-RNS) FHE implementation in OpenFHE against a 3-party secret-sharing SMPC implementation in MP-SPDZ. Both cryptographic pipelines execute equivalent financial transaction aggregation and multiplicative risk exposure workloads across datasets ranging from $N=100$ to $N=100,000$ transactions under Linux kernel traffic-controlled (\texttt{tc}/\texttt{netem}) latency profiles ($1\text{ ms}$ to $200\text{ ms}$) and bandwidth constraints ($1\text{ Gbps}$ down to $1\text{ Mbps}$). %%ABSTRACT_EMPIRICAL_SUMMARY%%
\end{abstract}

\begin{IEEEkeywords}
Fully Homomorphic Encryption, Secure Multi-Party Computation, OpenFHE, MP-SPDZ, Threshold BFV, Beaver Multiplication Triples, Secret Sharing, WAN Performance Crossover.
\end{IEEEkeywords}

\section{Introduction}
Cross-institutional financial analytics---such as systemic liquidity monitoring, syndicated fraud detection, and inter-bank risk exposure calculation---requires evaluating aggregate functions over private transaction ledgers held by mutually untrusting institutions. Traditional trusted third-party (TTP) data escrows introduce single points of breach and compliance liability. Cryptographic secure computation eliminates the trusted intermediary by allowing $M$ mutually untrusting parties $P_1, \dots, P_M$ with private inputs $D_1, \dots, D_M$ to evaluate a joint function $f(D_1, \dots, D_M)$ such that no participant learns anything beyond the authorized output.

Two mature software paradigms dominate privacy-preserving computation over arithmetic circuits:
\begin{enumerate}
    \item \textbf{Fully Homomorphic Encryption (FHE)}, which maps plaintexts into Ring-Learning-With-Errors (RLWE) polynomial quotient rings $R_q = \mathbb{Z}_q[X]/(X^n+1)$, supporting Single-Instruction Multiple-Data (SIMD) homomorphic addition and multiplication on encrypted ciphertexts without multi-round interaction during circuit evaluation~\cite{gentry2009,bfv2012,openfhe2022}.
    \item \textbf{Secure Multi-Party Computation (SMPC)}, which splits private inputs into linear secret shares over a finite field $\mathbb{F}_p$ or ring $\mathbb{Z}_{2^k}$, evaluating additions locally without communication and evaluating multiplications interactively via precomputed Beaver multiplication triples $([a], [b], [c])$~\cite{shamir1979,beaver1991,mpspdz2020}.
\end{enumerate}

While theoretical asymptotic complexities suggest that SMPC excels in low-latency Local Area Networks (LANs) and FHE excels in high-latency Wide Area Networks (WANs), real-world architectural selection requires precise quantitative crossover thresholds that account for SIMD batch packing, threshold key generation, Oblivious Transfer (OT) triple generation, ciphertext serialization expansion, and kernel-level TCP transport behavior.

\section{Problem Statement}
Consider three financial institutions---Participant A (\emph{Retail Banking}), Participant B (\emph{Commercial Treasury}), and Participant C (\emph{Cross-Border Settlement})---each holding a private partition $D_i = \{(x_{i,j}, w_{i,j})\}_{j=1}^{N_i}$ of a total dataset of $N = N_A + N_B + N_C$ financial transactions, where $x_{i,j} \in \mathbb{Z}^+$ denotes the transaction amount and $w_{i,j} \in \mathbb{Z}^+$ denotes a confidential counterparty risk score. The institutions must jointly execute two canonical workloads without disclosing individual records or intermediate private partitions:
\begin{itemize}
    \item \textbf{Workload 1 (Secure Aggregation):} Compute the exact global sum $\text{SUM} = \sum_{i=1}^3 \sum_{j=1}^{N_i} x_{i,j}$, total count $\text{COUNT} = N$, and exact arithmetic mean $\text{AVERAGE} = \text{SUM} / \text{COUNT}$ (where floating-point division is executed post-decryption/reconstruction on the authorized integer sum and count to prevent fixed-point rounding divergence).
    \item \textbf{Workload 2 (Multiplicative Risk Exposure):} Compute the inner-product weighted risk total $\text{RISK\_TOTAL} = \sum_{i=1}^3 \sum_{j=1}^{N_i} (x_{i,j} \cdot w_{i,j})$, where both $x_{i,j}$ and $w_{i,j}$ are protected private inputs requiring ciphertext-ciphertext homomorphic multiplication in FHE and Beaver triple interactive multiplication in SMPC.
\end{itemize}
The research problem is to experimentally measure and model how dataset cardinality $N$, one-way network propagation delay $L \in [1, 200]\text{ ms}$, and link capacity $B \in [1, 1000]\text{ Mbps}$ govern end-to-end execution latency, network traffic volume, memory footprint, and performance crossover boundaries between FHE and SMPC.

\section{Background: FHE and SMPC}
\subsection{Lattice-Based FHE and Threshold BFV-RNS}
In the Brakerski/Fan-Vercauteren (BFV) scheme~\cite{bfv2012}, plaintexts reside in the polynomial ring $R_p = \mathbb{Z}_p[X]/(X^n+1)$ for a power-of-two ring dimension $n$ and prime plaintext modulus $p \equiv 1 \pmod{2n}$, while ciphertexts consist of polynomial pairs $\mathbf{ct} = (c_0, c_1) \in R_q^2$ modulo a large composite ciphertext modulus $Q = \prod_{k=1}^K q_k$ represented via the Residue Number System (RNS). Under the Chinese Remainder Theorem (CRT), $R_p$ factors into $n$ independent slots, enabling SIMD batching of up to $n$ integers per ciphertext.

Encryption injects a small Gaussian error polynomial $e \in R_q$ ($\sigma = 3.2$). Homomorphic addition $\text{EvalAdd}(\mathbf{ct}_1, \mathbf{ct}_2)$ adds noise linearly, whereas homomorphic multiplication $\text{EvalMult}(\mathbf{ct}_1, \mathbf{ct}_2)$ forms a degree-2 tensor product $(c_0 c'_0, c_0 c'_1 + c_1 c'_0, c_1 c'_1)$ scaled by $\lfloor p/Q \rfloor$, followed by relinearization back to two ring elements using an evaluation key $\text{evk}$. Correct decryption holds only while $\|e_{\text{final}}\|_\infty < Q / (2p)$. To support multiple mutually untrusting data owners without a single secret-key holder, Multiparty Threshold FHE~\cite{mouchet2021} additively distributes the secret key $s = s_A + s_B + s_C$, constructs a collective public key $pk_{ABC}$ and collective relinearization key $\text{evk}_{ABC}$, and requires all three parties to compute partial decryptions padded with smudging/flooding noise (\texttt{NOISE\_FLOODING\_MULTIPARTY}) before fusion.

\subsection{Secret-Sharing SMPC and Beaver Multiplication Triples}
In linear secret-sharing SMPC over a prime field $\mathbb{F}_p$, a secret $x \in \mathbb{F}_p$ is split across $M=3$ parties either additively ($[x] = (x_1, x_2, x_3)$ with $\sum_{i=1}^3 x_i \equiv x \pmod p$, secure against a dishonest majority $t \le 2$) or via Shamir degree-$t$ polynomials ($f(0)=x$, secure against an honest majority $t < M/2 \implies t=1$)~\cite{shamir1979}. Any linear combination $\alpha [x] + \beta [y]$ is evaluated locally with zero communication.

Multiplying two secret-shared values $[x]$ and $[y]$ uses Beaver's randomization technique~\cite{beaver1991}. During an input-independent preprocessing phase, the parties generate fresh random shared triples $([a], [b], [c])$ satisfying $c = a \cdot b \pmod p$ (via Oblivious Transfer in dishonest-majority protocols~\cite{mascot2016}). In the online phase, parties locally compute masked shares $[d] = [x] - [a]$ and $[e] = [y] - [b]$, reconstruct $d$ and $e$ in one interactive broadcast round, and locally evaluate:
\begin{equation}
[x \cdot y] = [c] + d \cdot [b] + e \cdot [a] + d \cdot e \pmod p.
\end{equation}
Because $a$ and $b$ act as one-time pads masking $x$ and $y$, reusing a Beaver triple leaks linear correlations between private inputs and is strictly prohibited.

\section{Related Work}
Prior comparative studies of homomorphic encryption and multi-party computation have examined specialized domains such as federated logistic regression, neural network inference, and genomic association~\cite{openfhe2022,mpspdz2020,mouchet2021}. Al Badawi et al.~\cite{openfhe2022} introduced OpenFHE, unifying BGV, BFV, CKKS, and threshold multiparty extensions in C++ with Python bindings. Keller~\cite{mpspdz2020} developed MP-SPDZ, providing a unified virtual machine for over 30 honest- and dishonest-majority MPC protocols. However, existing comparisons frequently compare single-key FHE against multi-party MPC without aligning corruption thresholds, or estimate network transfer analytically rather than measuring kernel interface traffic under controlled WAN latency and bandwidth shaping. Our work addresses both gaps by aligning 3-party dishonest-majority threshold BFV against 3-party dishonest-majority OT-Beaver SMPC over real TCP sockets shaped by Linux \texttt{tc}/\texttt{netem}.

\section{System Architecture}
The nine-module experimental architecture implemented in this study consists of:
\begin{enumerate}
    \item \textbf{Dataset Generator (\texttt{dataset/}):} Produces deterministic, seeded ($\text{seed}=42$) synthetic transaction partitions for Participants A, B, and C, along with exact 64-bit integer reference results.
    \item \textbf{FHE Engine (\texttt{fhe/}):} Implements OpenFHE v1.6.1.0 context management, 3-party collective key generation, SIMD packing, binary serialization over TCP sockets, homomorphic evaluation, threshold noise-flooding decryption, and isolated CKKS bootstrapping benchmarks.
    \item \textbf{SMPC Engine (\texttt{smpc/}):} Compiles vectorized MP-SPDZ v0.4.4 \texttt{.mpc} programs and spawns three concurrent OS processes (\texttt{P0}, \texttt{P1}, \texttt{P2}) communicating over OpenSSL TLS channels with explicit four-phase timers.
    \item \textbf{Network Simulation Engine (\texttt{networking/}):} Applies Linux \texttt{tc}/\texttt{netem} qdisc rules for latency ($1$--$200\text{ ms}$) and bandwidth ($1$--$1000\text{ Mbps}$), validates live TCP RTT prior to each trial, and guarantees teardown cleanup.
    \item \textbf{Performance \& Traffic Monitor:} Samples process tree peak RSS memory and reads Linux kernel interface counters from \texttt{/proc/net/dev} (reporting \texttt{tx\_bytes} delta to prevent \texttt{tx+rx} double-counting on shared loopback/bridge interfaces).
    \item \textbf{Benchmark Orchestrator (\texttt{orchestrator/}):} Executes warm-up and measured repetitions, enforces cross-engine numerical equivalence, and exports raw and summarized CSV/JSON telemetry.
    \item \textbf{Statistical \& Crossover Analyzer (\texttt{analysis/}):} Performs Welch's $t$-tests, Mann-Whitney $U$ tests, Monte Carlo crossover interpolation, and decision matrix generation.
\end{enumerate}

\section{Experimental Methodology}
%%EXPERIMENT_STATUS_NOTE%%
To ensure fair comparison under aligned threat models, the primary experimental matrix compares 3-party Threshold BFV-RNS (\texttt{threshold\_bfv\_128}) against 3-party Dishonest-Majority Additive Secret Sharing with OT-based Beaver triples (\texttt{semi\_dishonest\_majority\_3p} via \texttt{semi-party.x}). Both primary protocols tolerate up to $t \le 2$ out of $3$ colluding semi-honest participants. Secondary baseline runs evaluate single-key BFV (\texttt{single\_key\_bfv\_128}) and honest-majority Shamir/Malicious-Shamir SMPC ($t \le 1$ of $3$), which are explicitly labeled as having non-equivalent trust assumptions.

Each configuration executes warm-up runs followed by repeated measured trials. We report parametric (mean, standard deviation, 95\% Student-$t$ confidence intervals) and non-parametric (median, min, max) metrics, separating setup/key-generation, offline preprocessing, input encryption/sharing, online evaluation, and output decryption/reconstruction.

\section{Cryptographic Implementation}
Table~\ref{tab:crypto_params} details the verified cryptographic parameters of both implementations. In OpenFHE BFV-RNS, we select the 29-bit NTT-friendly prime plaintext modulus $p = 536,903,681 \equiv 1 \pmod{65536}$ with SIMD batch size $4096$. Because transaction amounts lie in $[10, 200]$ and risk scores in $[1, 10]$, even the worst-case 100,000-transaction inner product ($2 \times 10^8$) remains strictly below the signed modulus bound $\lfloor p/2 \rfloor = 268,451,840$, guaranteeing exact integer arithmetic with zero modular wrap-around.

\begin{table}[htbp]
\caption{Cryptographic Parameter Specifications (FHE vs. SMPC)}
\label{tab:crypto_params}
\centering
\scriptsize
\begin{tabular}{@{}lll@{}}
\toprule
\textbf{Parameter} & \textbf{OpenFHE (Primary FHE)} & \textbf{MP-SPDZ (Primary SMPC)} \\
\midrule
%%CRYPTO_PARAM_TABLE_ROWS%%
\bottomrule
\end{tabular}
\end{table}

\section{Network Simulation Setup}
Network emulation is implemented via Linux kernel traffic control (\texttt{tc qdisc replace dev <iface> root netem delay <D>ms rate <R>mbit}). Five one-way delay profiles are evaluated: Local LAN ($1\text{ ms}$, $\text{RTT} \approx 2\text{ ms}$), Fast WAN ($10\text{ ms}$, $\text{RTT} \approx 20\text{ ms}$), Moderate WAN ($50\text{ ms}$, $\text{RTT} \approx 100\text{ ms}$), Slow WAN ($100\text{ ms}$, $\text{RTT} \approx 200\text{ ms}$), and High Latency WAN ($200\text{ ms}$, $\text{RTT} \approx 400\text{ ms}$). Four bandwidth tiers are tested: High ($1\text{ Gbps}$), Medium ($100\text{ Mbps}$), Restricted ($10\text{ Mbps}$), and Very Restricted ($1\text{ Mbps}$). Before every trial, a live \texttt{TCP\_NODELAY} socket probe measures actual round-trip time separately from configured one-way delay.

\section{Experimental Results}
All executed FHE and SMPC trials achieved 100\% numerical agreement with the 64-bit plaintext reference calculations across both workloads.

\subsection{Dataset Size Scaling}
Fig.~\ref{fig:graph1} and Table~\ref{tab:dataset_scaling} present end-to-end execution times across dataset sizes. Because OpenFHE packs $4096$ transactions into each RLWE ciphertext, increasing $N$ from $100$ to $1,000$ requires the exact same number of SIMD ciphertexts ($1$ ciphertext per participant for Workload 1; $2$ per participant for Workload 2), keeping FHE execution time virtually flat up to the ring batch boundary. Even at $N=100,000$ ($33,334$ records per organization, or $9$ packed ciphertexts per party), FHE requires only $27$ homomorphic multiplications and $26$ additions. Conversely, SMPC scales linearly with the number of Beaver triples $N$ generated during OT preprocessing and opened during online evaluation.

\begin{figure*}[htbp]
\centering
\includegraphics[width=0.92\textwidth]{../graphs/graph1_execution_time_vs_dataset_size.pdf}
\caption{Execution time vs. dataset size ($N$) for OpenFHE 3-Party Threshold BFV and MP-SPDZ 3-Party SMPC under Local LAN and WAN profiles.}
\label{fig:graph1}
\end{figure*}

\begin{table*}[htbp]
\caption{Experimental Performance Across Dataset Sizes (Local LAN: 1 ms Delay, 1 Gbps Bandwidth)}
\label{tab:dataset_scaling}
\centering
\scriptsize
\begin{tabular}{@{}llrrrrrrr@{}}
\toprule
\textbf{Workload} & \textbf{Dataset ($N$)} & \textbf{FHE Mean$\pm$SD (s)} & \textbf{SMPC Mean$\pm$SD (s)} & \textbf{FHE Traffic (MB)} & \textbf{SMPC Traffic (MB)} & \textbf{FHE Mem (MB)} & \textbf{SMPC Mem (MB)} & \textbf{Faster (Speedup)} \\
\midrule
%%DATASET_SCALING_TABLE_ROWS%%
\bottomrule
\end{tabular}
\end{table*}

\subsection{Network Latency and Bandwidth Sensitivity}
Figs.~\ref{fig:graph2}, \ref{fig:graph3}, and \ref{fig:graph4} illustrate the decoupled effects of network latency and bandwidth, summarized quantitatively in Table~\ref{tab:network_sweeps}. As shown in Fig.~\ref{fig:graph2}, aggregate network traffic volume is invariant to propagation latency for both technologies, but differs substantially in composition: Threshold FHE exchanges fixed-size RLWE public keys, collective relinearization keys, and packed ciphertexts in $O(1)$ communication rounds, whereas SMPC exchanges compact field shares and OT payloads over $14$--$30+$ interactive rounds.

\begin{figure*}[htbp]
\centering
\includegraphics[width=0.90\textwidth]{../graphs/graph2_network_traffic_vs_latency.pdf}
\caption{Aggregate measured network traffic (MB) vs. configured one-way network latency (ms).}
\label{fig:graph2}
\end{figure*}

\begin{figure*}[htbp]
\centering
\includegraphics[width=0.90\textwidth]{../graphs/graph3_execution_time_vs_latency.pdf}
\caption{End-to-end execution time (s) vs. one-way network latency (ms) at 1 Gbps bandwidth, highlighting observed latency crossover thresholds.}
\label{fig:graph3}
\end{figure*}

\begin{figure*}[htbp]
\centering
\includegraphics[width=0.90\textwidth]{../graphs/graph4_execution_time_vs_bandwidth.pdf}
\caption{End-to-end execution time (s) vs. available network bandwidth (Mbps), demonstrating FHE sensitivity to severe bandwidth throttling (1 Mbps).}
\label{fig:graph4}
\end{figure*}

\begin{table*}[htbp]
\caption{Experimental Results Across Independent Network Latency and Bandwidth Sweeps ($N=1,000$, Workload 2: Multiplicative Risk)}
\label{tab:network_sweeps}
\centering
\scriptsize
\begin{tabular}{@{}llrrrrrr@{}}
\toprule
\textbf{Sweep Dimension} & \textbf{Network Condition} & \textbf{Measured RTT (ms)} & \textbf{FHE Mean$\pm$SD (s)} & \textbf{SMPC Mean$\pm$SD (s)} & \textbf{$\Delta T = T_{\text{FHE}}-T_{\text{SMPC}}$ (s)} & \textbf{Welch $p$-value} & \textbf{Winner (Speedup)} \\
\midrule
%%NETWORK_SWEEP_TABLE_ROWS%%
\bottomrule
\end{tabular}
\end{table*}

\section{Comparative Analysis: Storage, Phases, and Noise}
Fig.~\ref{fig:graph5} compares cryptographic payload expansion against raw 64-bit integer plaintexts, and Fig.~\ref{fig:graph6} decomposes execution time into constituent phases. In SMPC, each 8-byte secret integer is represented by one 16-byte prime-field share per party across 3 parties ($48\text{ bytes}$ total), yielding a constant $6.0\times$ share storage expansion ratio regardless of $N$. In contrast, an OpenFHE BFV-RNS ciphertext occupies a fixed polynomial ring footprint regardless of how many of its $4096$ SIMD slots are populated; thus, FHE expansion ratio drops hyperbolically as $N$ grows and amortizes the ring dimension.

%%NOISE_AND_BOOTSTRAPPING_PARAGRAPH%%

\begin{figure*}[htbp]
\centering
\includegraphics[width=0.90\textwidth]{../graphs/graph5_storage_expansion_comparison.pdf}
\caption{Cryptographic payload size (KB) and storage expansion ratio across dataset sizes for OpenFHE SIMD ciphertexts and MP-SPDZ 3-party shares.}
\label{fig:graph5}
\end{figure*}

\begin{figure*}[htbp]
\centering
\includegraphics[width=0.88\textwidth]{../graphs/graph6_computation_phase_breakdown.pdf}
\caption{Stacked execution phase breakdown comparing Setup/KeyGen, Preprocessing/Beaver Triples, Encryption/Sharing, Online Evaluation, and Decryption/Reconstruction.}
\label{fig:graph6}
\end{figure*}

\section{Crossover Threshold Analysis}
Table~\ref{tab:crossovers} reports the experimentally observed and piecewise-linearly interpolated crossover thresholds along with 95\% Monte Carlo uncertainty bounds. Following strict experimental methodology, measured bracketing intervals and interpolated thresholds are reported separately, and regimes without a sign change in $\Delta T = T_{\text{FHE}} - T_{\text{SMPC}}$ are explicitly labeled as having no crossover in the tested range. Table~\ref{tab:decision_matrix} synthesizes the experimentally derived decision criteria.

\begin{table*}[htbp]
\caption{Experimentally Observed and Interpolated Performance Crossover Thresholds}
\label{tab:crossovers}
\centering
\scriptsize
\begin{tabular}{@{}llllll@{}}
\toprule
\textbf{Axis} & \textbf{Workload} & \textbf{Fixed Conditions} & \textbf{Transition} & \textbf{Measured Bracket} & \textbf{Interpolated Threshold (95\% CI)} \\
\midrule
%%CROSSOVER_TABLE_ROWS%%
\bottomrule
\end{tabular}
\end{table*}

\begin{table*}[htbp]
\caption{Experimentally Derived Decision Matrix for Private Collaborative Computing}
\label{tab:decision_matrix}
\centering
\scriptsize
\begin{tabular}{@{}p{0.28\textwidth}p{0.14\textwidth}p{0.52\textwidth}@{}}
\toprule
\textbf{Operational Conditions} & \textbf{Recommended Approach} & \textbf{Experimental Evidence (Mean$\pm$SD, Traffic, Significance)} \\
\midrule
%%DECISION_MATRIX_ROWS%%
\bottomrule
\end{tabular}
\end{table*}

\section{Discussion}
The empirical findings demonstrate a fundamental architectural duality between computation-bound lattice homomorphic encryption and communication-round-bound secret-sharing MPC:
\begin{enumerate}
    \item \textbf{Why Latency Dominates SMPC:} Because OT-based Beaver triple generation and online Beaver openings require $\sim 20$--$30$ sequential synchronous rounds, every increase of $\Delta L$ in one-way propagation delay adds $\approx 2 \cdot R \cdot \Delta L$ to SMPC wall-clock execution time. Once one-way WAN delay exceeds the latency crossover threshold, SMPC's round-trip stall time eclipses FHE's local polynomial evaluation cost.
    \item \textbf{Why Bandwidth Throttling Penalizes Threshold FHE:} Threshold BFV requires exchanging multi-megabyte collective relinearization keys and RLWE ciphertexts. When link capacity drops to $1\text{ Mbps}$ (\emph{Very Restricted Bandwidth}), serializing and transmitting FHE key and ciphertext payloads over TCP introduces multi-second serialization transfer delays, whereas SMPC's compact field shares for small-to-medium datasets traverse a $1\text{ Mbps}$ link in milliseconds.
    \item \textbf{Why SIMD Batching Favors FHE at Scale:} For vectorizable inner-product workloads ($N \ge 10,000$), SIMD slot packing allows a single OpenFHE \texttt{EvalMult} operation to evaluate $4096$ products simultaneously, whereas SMPC must generate and open $N$ distinct Beaver triples.
\end{enumerate}

\section{Limitations and Threat Model Differences}
First, our primary benchmark operates under the semi-honest (passive) adversary model with dishonest majority ($t \le 2$ of $3$). Upgrading SMPC to active malicious security with abort (\texttt{malicious-shamir-party.x} or \texttt{mascot-party.x}) introduces triple sacrifice or MAC verification overhead, whereas upgrading threshold FHE to active malicious security requires zero-knowledge proofs of plaintext knowledge ($\Sigma$-protocols) during encryption and partial decryption. Second, our network emulation shapes a shared container interface (\texttt{lo}/\texttt{eth0}) with deterministic \texttt{netem} queues; real-world transcontinental BGP routing exhibits bursty cross-traffic and TCP congestion window dynamics beyond steady-state \texttt{netem} shaping.

\subsection{Reproducibility via Docker}
To reproduce the entire benchmark suite, statistical crossover analysis, graphs, and this IEEE PDF report inside an isolated Ubuntu 24.04 container with \texttt{NET\_ADMIN} capability for \texttt{tc}/\texttt{netem}:
\begin{verbatim}
docker compose run --build --rm benchmark
\end{verbatim}
For a rapid verification run across the quick matrix:
\begin{verbatim}
docker compose run --rm benchmark \
  python3 -m orchestrator.runner --mode quick
\end{verbatim}

\section{Conclusion}
We have designed, implemented, and experimentally validated a reproducible comparative benchmark of 3-party Threshold BFV Fully Homomorphic Encryption (OpenFHE) and 3-party Secret-Sharing Multi-Party Computation (MP-SPDZ) across synthetic financial aggregation and multiplicative risk workloads. Neither cryptographic paradigm is universally faster: SMPC provides superior performance for small-to-moderate datasets on low-latency links or severely bandwidth-constrained channels, whereas SIMD-packed Threshold FHE outperforms SMPC in high-latency WAN environments and large-scale batch multiplicative workloads where $O(1)$ round complexity and SIMD amortization outweigh lattice polynomial overhead.

\begin{thebibliography}{10}
\bibitem{gentry2009} C. Gentry, ``Fully homomorphic encryption using ideal lattices,'' in \emph{Proc. 41st Annu. ACM Symp. Theory Comput. (STOC)}, 2009, pp. 169--178.
\bibitem{bfv2012} J. Fan and F. Vercauteren, ``Somewhat practical fully homomorphic encryption,'' \emph{IACR Cryptol. ePrint Arch.}, Rep. 2012/144, 2012.
\bibitem{openfhe2022} A. Al Badawi \emph{et al.}, ``OpenFHE: Open-source fully homomorphic encryption library,'' in \emph{Proc. 10th Workshop Encrypted Comput. Appl. Homomorphic Cryptogr. (WAHC)}, ACM, 2022, pp. 53--63.
\bibitem{mouchet2021} C. Mouchet, J. Troncoso-Pastoriza, J.-P. Bossuat, and J.-P. Hubaux, ``Multiparty homomorphic encryption from ring-learning-with-errors,'' \emph{Proc. Privacy Enhancing Technol. (PoPETs)}, vol. 2021, no. 4, pp. 291--311, 2021.
\bibitem{shamir1979} A. Shamir, ``How to share a secret,'' \emph{Commun. ACM}, vol. 22, no. 11, pp. 612--613, Nov. 1979.
\bibitem{beaver1991} D. Beaver, ``Efficient multiparty protocols using circuit randomization,'' in \emph{Advances in Cryptology (CRYPTO)}, Springer, 1991, pp. 420--432.
\bibitem{mpspdz2020} M. Keller, ``MP-SPDZ: A versatile framework for multi-party computation,'' in \emph{Proc. ACM SIGSAC Conf. Comput. Commun. Security (CCS)}, 2020, pp. 1575--1590.
\bibitem{spdz2012} I. Damg{\aa}rd, V. Pastro, N. Smart, and S. Zakarias, ``Multiparty computation from somewhat homomorphic encryption,'' in \emph{Advances in Cryptology (CRYPTO)}, Springer, 2012, pp. 643--662.
\bibitem{mascot2016} M. Keller, E. Orsini, and P. Scholl, ``MASCOT: Faster malicious arithmetic secure computation with oblivious transfer,'' in \emph{Proc. ACM SIGSAC Conf. Comput. Commun. Security (CCS)}, 2016, pp. 830--842.
\bibitem{lindell2017} Y. Lindell and A. Nof, ``A framework for constructing fast MPC over arithmetic circuits with malicious adversaries and an honest-majority,'' in \emph{Proc. ACM SIGSAC Conf. Comput. Commun. Security (CCS)}, 2017, pp. 259--276.
\end{thebibliography}

\end{document}
